\documentclass[../main.tex]{subfiles}
\begin{document}

% Эх: report/2026-dadlaga_tulvlguu_last.docx (2026.06.15).
% Төлөвлөсөн ажлын англи эх, 12 биелэлт, удирдагчийн тайлбарыг хэвээр хадгалав.
% BullMQ/SSE нь анхны төлөвлөгөө; хэрэгжилтийн зөрүүг дүгнэлтэд тайлбарласан.
\begin{center}
\textbf{RAG CHATBOT}
\end{center}
\begin{flushright}
2026 он 06 сар 15 өдөр
\end{flushright}

\begingroup
\fontsize{10}{12}\selectfont
\renewcommand{\arraystretch}{1.25}
\setlength{\tabcolsep}{4pt}
\begin{longtable}{|>{\centering\arraybackslash}p{0.5cm}
                  |>{\raggedright\arraybackslash}p{5.3cm}
                  |>{\centering\arraybackslash}p{1.7cm}
                  |>{\centering\arraybackslash}p{1.3cm}
                  |>{\raggedright\arraybackslash}p{\dimexpr\linewidth-8.8cm-10\tabcolsep-6\arrayrulewidth\relax}|}
\caption{Үйлдвэрлэлийн дадлагын ажлын төлөвлөгөө, биелэлт}
\label{tab:internship-plan}
\\ \hline
\textbf{№} & \textbf{Гүйцэтгэх ажил} & \textbf{Хугацаа} &
\textbf{Биелэлт} & \textbf{Дадлагын удирдагчийн үнэлгээ}
\\ \hline
\endfirsthead

\multicolumn{5}{r}{\small Хүснэгт~\thetable-ын үргэлжлэл}\\[0.4em]
\hline
\textbf{№} & \textbf{Гүйцэтгэх ажил} & \textbf{Хугацаа} &
\textbf{Биелэлт} & \textbf{Дадлагын удирдагчийн үнэлгээ}
\\ \hline
\endhead
\endfoot
\endlastfoot

1 &
Setup: monorepo, Postgres + pgvector, schema (documents, chunks, conversations, messages), HNSW index &
06.15 &
сайн &
Monorepo орчин, PostgreSQL + pgvector өгөгдлийн сан, шаардлагатай хүснэгтүүд болон HNSW индексийг тохируулж хэрэгжүүлсэн.
\\ \hline

2 &
Dual ingestion paths in-memory: (a) file upload $\rightarrow$ buffer $\rightarrow$ extract text, (b) text paste $\rightarrow$ straight to chunking &
06.16 &
сайн &
PDF файл байршуулж текст ялгах болон текст шууд оруулах хоёр урсгалыг хэрэгжүүлсэн.
\\ \hline

3 &
Chunk ($\sim$500--800 tokens, 15\% overlap) $\rightarrow$ batch embed $\rightarrow$ store vectors in pgvector, via BullMQ worker (processing $\rightarrow$ ready) &
06.17 &
сайн &
Текстийг 500--800 токеноор хэсэглэн, embedding үүсгэж pgvector-д хадгалах ажиллагааг хэрэгжүүлсэн.
\\ \hline

4 &
Retrieval: query embedding $\rightarrow$ cosine top-k (5--8) + similarity threshold; test quality in isolation &
06.18 &
сайн &
Асуултын embedding үүсгэн cosine similarity болон top-k хайлтаар холбогдох хэсгүүдийг сонгох функцийг хэрэгжүүлж туршсан.
\\ \hline

5 &
Generation: “answer only from context” prompt + retrieved chunks, streamed via SSE &
06.19 &
сайн &
Олдсон контекстэд тулгуурлан хариу үүсгэх зааварчилгааг хэрэгжүүлсэн. SSE streaming-г цаашдын хөгжүүлэлтэд үлдээсэн.
\\ \hline

6 &
Chat UI: streaming chat, upload + paste with status, conversation list, source display &
06.22-23 &
сайн &
Вэб чат интерфэйс, файл/текст оруулах, боловсруулалтын төлөв, харилцан ярианы жагсаалт, эх сурвалж харах боломжийг хэрэгжүүлсэн.
\\ \hline

\pagebreak

7 &
Conversation memory: persist messages, re-retrieve per question &
06.24 &
сайн &
Харилцан яриа болон мессежүүдийг өгөгдлийн санд хадгалж, асуулт бүрд дахин хайлт хийх урсгалыг хэрэгжүүлсэн.
\\ \hline

8 &
Edge cases: empty retrieval, large files, multi-doc, rate limiting, token budget &
06.25 &
сайн &
Хоосон хайлт, том хэмжээний текст, олон баримт, API хүсэлтийн хязгаарлалт, токены төсөвтэй холбоотой нөхцөлүүдийг шалгаж, алдаа боловсруулах логик нэмсэн.
\\ \hline

9 &
Citations: return source snippets (filename, page) with answers &
06.26 &
сайн &
Хариултад ашигласан эх сурвалжийн файлын нэр, хуудас болон холбогдох хэсгийг буцаах боломжийг хэрэгжүүлсэн.
\\ \hline

10 &
Auth \& multi-tenancy: scope all docs/queries to user\_id, keys server-side only, quotas &
06.29-30 &
сайн &
Хэрэглэгчийн бүртгэл, нэвтрэлт, session шалгалт болон өгөгдлийг user\_id-аар тусгаарлах хамгаалалтыг хэрэгжүүлсэн.
\\ \hline

11 &
Testing: Q\&A eval set against known docs, load-test ingestion &
07.01 &
сайн &
Үндсэн урсгалын smoke test, typecheck, build болон баримт боловсруулах ачааллын туршилтуудыг хийсэн.
\\ \hline

12 &
Deploy: Dockerize, monitoring/logging &
07.02 &
сайн &
Docker Compose ашиглан локал байршуулалтын орчныг бэлтгэж, лог болон хяналтын үндсэн тохиргоог хийсэн.
\\ \hline

\end{longtable}
\endgroup

\vspace{1em}
\noindent Баталсан: Албан байгууллагын дадлага удирдагч: \dotfill

\vspace{1em}
\noindent Төлөвлөгөө боловсруулсан: Оюутан \dotfill

\end{document}
